% Public website edition. Personal telephone and postal address are intentionally omitted.
% Keep PNNL-DS-October.tex for job applications.
\documentclass{resume} % Use the custom resume.cls style

\usepackage{comment}
\usepackage[left=0.8 in,top=0.6in,right=0.8 in,bottom=0.8in]{geometry} % Document margins
\newcommand{\tab}[1]{\hspace{.2667\textwidth}\rlap{#1}} 
\newcommand{\itab}[1]{\hspace{0em}\rlap{#1}}
\usepackage{float}
\usepackage{fancyhdr}
\usepackage{tabularx}
\usepackage{array}
\pagestyle{fancy}
\fancyhf{}
\rfoot{\thepage}

\begin{document}
\vspace{-0.5in}
\begin{center}
    \textbf{\Huge \scshape Siddhartha Shankar Das} \\
    \vspace{1pt}
    Email: \href{mailto:siddhartha.das@pnnl.gov}{\underline{siddhartha.das@pnnl.gov}} /
    \href{mailto:siddhartha047@gmail.com}{\underline{siddhartha047@gmail.com}} \\
    \vspace{1pt}
    \href{https://scholar.google.com/citations?user=kTAFl2yYe6QC\&hl=en}{\underline{Google Scholar}} $|$
    \href{https://www.linkedin.com/in/siddhartha-shankar-das-25993b110/}{\underline{LinkedIn}} 
    $|$
    \href{https://sites.google.com/site/siddhartha047/}{\underline{Website}} 
    $|$
    \href{https://www.pnnl.gov/people/siddhartha-shankar-das}{\underline{PNNL Staff Page}}
\end{center}

\begin{rSection}{Interests}
I am a Post-Doctorate Researcher at Pacific Northwest National Laboratory (PNNL) in the \textbf{Language Intelligence} group within the \textbf{Data Science and Machine Intelligence} department. My research focuses on artificial intelligence and machine learning for scientific and applied domains, with interests spanning graph neural networks, graph sparsification, natural language processing, cybersecurity, and scalable AI systems. I develop learning-based methods for system design and optimization, including hardware/software partitioning and efficient learning on large graphs. I also build LLM-based, multimodal, and agentic AI systems for document intelligence, retrieval-augmented generation, georeferencing, and interactive decision support. My work has been published in top-tier venues and has received best paper awards.
\end{rSection}

\begin{rSection}{Education}

{\bf Purdue University, West Lafayette, IN, USA}\\
{Ph.D. in Computer Science}\hfill {Aug 2018 -- May 2025}\\
Advisor: \href{https://www.cs.purdue.edu/homes/apothen/}{Alex Pothen}\\
\textbf{Thesis:} Scaling Graph Neural Networks for Large Graphs and Analyzing Common Vulnerabilities via Large Language Models [\href{https://hammer.purdue.edu/articles/thesis/Scaling_Graph_Neural_Networks_for_Large_Graphs_and_Analyzing_Common_Vulnerabilities_via_Large_Language_Models/28856864}{pdf}].\\
\textbf{Topics:} Machine Learning on Graphs, Graph Neural Networks, Cybersecurity, Large Language Models \\

{\bf Bangladesh University of Engineering and Technology (BUET), Dhaka, Bangladesh}\\
{M.Sc. in Computer Science and Engineering}\hfill {Sep 2017}\\
GPA: 4.00/4.00\\
\textbf{Thesis:} Evolutionary Algorithm Using Adaptive Fuzzy Dominance and Reference Point for Many-Objective Optimization ($F$-DEA) [\href{https://drive.google.com/open?id=1-hbCL9_WCNSZu0qSFklLzoJqylzHG_2k}{pdf}].\\
Advisor: \href{https://cse.buet.ac.bd/faculty/facdetail.php?id=mdmonirulislam}{Md. Monirul Islam}\\
\textbf{Topics:} Multi-objective Optimization, Evolutionary Algorithms \\

{\bf Bangladesh University of Engineering and Technology (BUET), Dhaka, Bangladesh}\\
{B.Sc. in Computer Science and Engineering}\hfill {Jul 2014}\\
GPA: 3.88/4.00, Major: 3.99/4.00\\
\textbf{Thesis:} A Many-objective Evolutionary Approach Using Fuzzy Dominance with Bidirectional Bias [\href{https://drive.google.com/open?id=0B-SybtAwa8vccEI0UWkybGdWVzg}{pdf}].\\
Advisor: \href{https://cse.buet.ac.bd/faculty/facdetail.php?id=mdmonirulislam}{Md. Monirul Islam}\\
\textbf{Topics:} Multi-objective Optimization, Evolutionary Algorithms

\end{rSection}

\begin{rSection}{Work Experience}
\textbf{Postdoctoral Researcher, Data Science \& Machine Intelligence} \hfill May 2025 -- Present\\
\href{https://www.pnnl.gov/}{Pacific Northwest National Laboratory (PNNL)}, Richland, WA, USA.\\
\begin{itemize}
\item Conduct research in artificial intelligence, machine learning, graph algorithms, and system design, with a focus on scalable Graph Neural Network (GNN)-based methods for optimization problems such as hardware/software partitioning, scheduling, graph partitioning, and graph sparsification.

\item Develop LLM-based, multimodal, and agentic AI systems for scientific and applied applications, including large-scale processing of environmental and regulatory documents. Current work includes metadata extraction, georeferencing, project summarization, retrieval-augmented generation, and interactive decision-support systems.

\item Contribute to interdisciplinary projects spanning cybersecurity, document intelligence, and scientific machine learning, translating foundational methods into scalable tools, datasets, deployable systems, and research publications.
\end{itemize}

\textbf{Graduate Research Assistant} \hfill Aug 2018 -- May 2025\\
Department of Computer Science, Purdue University, West Lafayette, IN, USA.\\
\begin{itemize}
\item Conducted research on scalable graph learning, graph sparsification, and machine learning on large graphs, along with LLM-based methods for software vulnerability analysis and cyber threat understanding.

\item Developed new sampling and sparsification methods for Graph Neural Networks, leading to publications in venues including KDD, DSAA, Big Data, and HST.

% \item Completed Ph.D. thesis: \textit{Scaling Graph Neural Networks for Large Graphs and Analyzing Common Vulnerabilities via Large Language Models.}
\end{itemize}

\textbf{Ph.D. Intern, Data Analytics} \hfill Summers 2019--2021\\
\href{https://www.pnnl.gov/}{Pacific Northwest National Laboratory (PNNL)}, Richland, WA, USA\\
Mentored by \href{https://www.pnnl.gov/people/mahantesh-halappanavar}{Mahantesh Halappanavar}.\\
\begin{itemize}
    \item Developed learning-based methods for classification and analysis of software vulnerabilities, including mapping CVEs to weaknesses (CWEs) and attack patterns (CAPECs).
    \item Used Large Language Models (LLMs) and Natural Language Processing (NLP) techniques, and scaled these methods for modern AI accelerators.
    \item Published four papers during the internship period, including work recognized with two best paper awards.
\end{itemize} 

\textbf{Assistant Professor, Department of Computer Science and Engineering} \hfill Dec 2017 -- Jul 2018\\
\href{https://www.buet.ac.bd/web/}{Bangladesh University of Engineering and Technology (BUET)}, Dhaka, Bangladesh\\
\textbf{Roles:} University teaching, project supervision, and office administration.\\
\textbf{Study Leave for Ph.D.} Jul 2018 -- Jul 2024

\textbf{Lecturer, Department of Computer Science and Engineering} \hfill Aug 2014 -- Dec 2017\\
\href{https://www.buet.ac.bd/web/}{Bangladesh University of Engineering and Technology (BUET)}, Dhaka, Bangladesh\\
\textbf{Roles:} University teaching.

\end{rSection} 

\begin{rSection}{Selected Publications (\textbf{Full list at } \href{https://scholar.google.com/citations?user=kTAFl2yYe6QC&hl=en}{Google Scholar})}

\textbf{Conference and Workshop Publications}
\begin{itemize}
    \item \textbf{Siddhartha Shankar Das}, N. Arafat, M. Rahman, S. M. Ferdous, A. Pothen, M. Halappanavar, and D. Rawat. ``SGS-GNN: A Supervised Graph Sparsifier for Graph Neural Networks." \textit{In Proceedings of the 32nd ACM SIGKDD Conference on Knowledge Discovery and Data Mining V.2 (KDD 2026), August 9–13, 2026, 12 pages.} \href{https://doi.org/10.1145/3770855.3818110}{[doi]}

    \item \textbf{Siddhartha Shankar Das}, Zakaria Mehrab, James Kotary, Rounak Meyur, S. M. Ferdous, Erdal Mutlu, and Mahantesh Halappanavar. “Diff-GNN: Differentiable Graph Neural Networks for Joint Hardware–Software Partitioning and Scheduling.” Accepted at HPEC 2026.
    
    \item James Kotary, Zakaria Mehrab, \textbf{Siddhartha Shankar Das}, Rounak Meyur, S. M. Ferdous, Erdal Mutlu, and Mahantesh Halappanavar. “Joint Optimization of HW/SW Partitioning and Scheduling with Deep Reinforcement Learning.” Accepted at HPEC 2026.
    
    \item Moqsadur Rahman, Monika Akbar, Joseph T. Aguayo, Mahmud Shahriar Hossain, Tina M. Ellis, \textbf{Siddhartha Shankar Das} et al. ``ATT\&CK to CVE: A Large-Scale Automated Knowledge Graph for Threat Intelligence." IEEE International Conference on Data Mining (ICDM) GTA'3 Workshop, 2025. \href{https://monowarhasan.info/papers/attack2cve_GTA3_25.pdf}{[paper]}
    
    \item \textbf{Siddhartha Shankar Das}, S. M. Ferdous, M. Halappanavar, E. Serra, and A. Pothen. ``AGS-GNN: Attribute-guided Sampling for Graph Neural Networks." In \textit{Proceedings of the 30th ACM SIGKDD Conference on Knowledge Discovery and Data Mining (KDD)}, 2024. \href{https://arxiv.org/abs/2405.15218}{[paper]}
    
    \item \textbf{Siddhartha Shankar Das}, E. Serra, M. Halappanavar, A. Pothen, and E. Al-Shaer. ``V2W-BERT: A Framework for Effective Hierarchical Multiclass Classification of Software Vulnerabilities." In \textit{2021 IEEE 8th International Conference on Data Science and Advanced Analytics (DSAA)}, pp. 1--12, 2021. \href{https://doi.org/10.1109/DSAA53316.2021.9564227}{[doi]} [\href{https://www.cs.purdue.edu/news/articles/2021/pothen_v2w-bert.html}{\textbf{Best Application Paper}}]

    \item \textbf{Siddhartha Shankar Das}, A. Dutta, S. Purohit, E. Serra, M. Halappanavar, and A. Pothen. ``Towards Automatic Mapping of Vulnerabilities to Attack Patterns Using Large Language Models." In \textit{2022 IEEE International Symposium on Technologies for Homeland Security (HST)}, 2022. \href{https://doi.org/10.1109/HST56032.2022.10025459}{[doi]} [\href{https://lh4.googleusercontent.com/RgzcFOh0gfXK4MTvMchVhGQ-X39kd13UFPGsZuHLkCcC9-CJNQ2dW_nfg0pIcuCSYbpbWlt94_z1dST95MYITwod_aW0tvTmpOLhRZqZs9eC73thuy5eA4YSDZLll9EK-w=w1280}{\textbf{Best Paper}}]

    \item \textbf{Siddhartha Shankar Das}, M. Halappanavar, A. Tumeo, E. Serra, A. Pothen, and E. Al-Shaer. ``VWC-BERT: Scaling Vulnerability--Weakness--Exploit Mapping on Modern AI Accelerators." In \textit{2022 IEEE International Conference on Big Data}, pp. 1224--1229, 2022. \href{https://doi.org/10.1109/BigData55660.2022.10020622}{[doi]}

    \item K. Panchal, \textbf{Siddhartha Shankar Das}, L. Torre, J. Miller, R. Rallo, and M. Halappanavar. ``Efficient Clustering of Software Vulnerabilities Using Self-Organizing Maps (SOM)." In \textit{2022 IEEE International Symposium on Technologies for Homeland Security (HST)}, pp. 1--7, 2022. \href{https://doi.org/10.1109/HST56032.2022.10025443}{[doi]}

    \item A. Rizvi, T. R. Toha, \textbf{Siddhartha Shankar Das}, Sriram Chellappan, A. B. M. Alim Al Islam. ``Many-objective Performance Enhancement in Computing Clusters." In \textit{2017 IEEE 36th International Performance Computing and Communications Conference (IPCCC)}, 2017. \href{https://doi.org/10.1109/PCCC.2017.8280491}{[doi]}

    \item A. Rizvi, T. R. Toha, \textbf{Siddhartha Shankar Das}, Sriram Chellappan, A. B. M. Alim Al Islam. ``Exploiting a Synergy between Greedy Approach and NSGA for Scheduling in Computing Clusters." In \textit{2018 5th International Conference on Networking, Systems and Security}, 2018. \href{https://doi.org/10.1109/NSysS.2018.8631376}{[doi]}
\end{itemize}

\textbf{Journal Articles, Book Chapters, and Reports}
\begin{itemize}
    \item Moqsadur Rahman, Aaron Sanchez, Krish Piryani, \textbf{Siddhartha Shankar Das}, Sai Munikoti, Luis de la Torre Quintana, Monowar Hasan, Joseph Aguayo, Monika Akbar, Shahriar Hossain, and Mahantesh Halappanavar. ``Generative AI for Advanced Cyber Defense." \textit{AI for Cybersecurity: Research and Practice}, John Wiley \& Sons, Inc., pp. 109--146, 2025. \href{https://onlinelibrary.wiley.com/doi/abs/10.1002/9781394293773.ch04}{[book chapter]}

    \item Sai Munikoti, Dan Nally, Sai Dileep Koneru, \textbf{Siddhartha Shankar Das}, et al. ``NEPATEC v2.0: Standardized Metadata and Text Corpus of National Environmental Policy Act Documents." Pacific Northwest National Laboratory, 2025. \href{https://www.pnnl.gov/publications/nepatec-v20-standardized-metadata-and-text-corpus-national-environmental-policy-act}{[report]}

    \item \textbf{Siddhartha Shankar Das}, Md. Monirul Islam, Naheed Anjum Arafat. ``Evolutionary Algorithm Using Adaptive Fuzzy Dominance and Reference Point for Many-Objective Optimization." \textit{Swarm and Evolutionary Computation}, 44:1092--1107, 2019. \href{https://doi.org/10.1016/j.swevo.2018.11.003}{[doi]}
\end{itemize}

\textbf{Preprints}
\begin{itemize}

    \item \textbf{Siddhartha Shankar Das}, S. K. Navuluru, S. M. Ferdous, R. A. Rossi, B. Coskunuzer, L. Tamil, E. Serra, A. Pothen, R. Rallo, and M. M. Halappanavar. ``Scaffold: Support Graph Theory Based Sparsification for Graph Neural Networks." \textit{arXiv:2609.31466}, 2026. \href{https://arxiv.org/abs/2609.31466}{[arxiv]}

    \item S. K. Navuluru, \textbf{Siddhartha Shankar Das}, F. Dernoncourt, S. M. Ferdous, R. A. Rossi, N. K. Ahmed, B. Coskunuzer, A. Pothen, L. Tamil, and M. M. Halappanavar. ``EDiS: Edge Disjoint Subgraph Sparsification Framework for Graph Neural Networks." \textit{arXiv:2610.09059}, 2026. \href{https://arxiv.org/abs/2610.09059}{[arxiv]}

    \item S. K. Navuluru, \textbf{Siddhartha Shankar Das}, B. Ni, H. Chen, Y. Wang, B. Coskunuzer, N. K. Ahmed, F. Dernoncourt, M. Halappanavar, T. Derr, R. A. Rossi, and L. Tamil. ``Unifying Graph Neural Networks Through a Common Layer Equation." \textit{arXiv:2608.16097}, 2026. \href{https://arxiv.org/abs/2608.16097}{[arxiv]}

\end{itemize}
\end{rSection}



\begin{rSection}{Ongoing and Past Projects}

\textbf{Graph Learning, System Design, and Optimization}
\begin{itemize}
    \item \textbf{Graph-Based Circuit Equivalence Checking:}
    Developing graph learning methods for circuit equivalence verification using GNN-based subcircuit identification, gate embeddings, and graph matching.

    \item \textbf{Benchmarking Graph Sparsification Methods and Applications:}
    Developing a comprehensive textbook and benchmarking framework covering graph sparsification, rewiring, augmentation, theoretical guarantees, and applications.

    \item \textbf{Graph Neural Network Architecture Generation:}
    Exploring automated GNN architecture generation through a unified framework for message passing, aggregation, and transformation.
\end{itemize}

\textbf{NEPA, Multimodal AI, and Agentic Systems}
\begin{itemize}
    \item \textbf{NEPATEC 2.0: Standardized NEPA Document Corpus:}
    Contributed to a publicly available NEPA dataset for document intelligence, released on \href{https://huggingface.co/datasets/PNNL/NEPATEC2.0}{Hugging Face}.

    \item \textbf{NEPATEC 2.5: Multimodal Metadata Extraction:}
    Developing multimodal LLM methods for NEPA metadata extraction, map-based geolocation, and project summarization.

    \item \textbf{PermitCE: AI-Assisted Chat for NEPA Documents:}
    Developing a retrieval-augmented generation system for semantic search, recommendations, summarization, and question answering.
\end{itemize}

\textbf{Cybersecurity and Applied Machine Learning}
\begin{itemize}
    \item \textbf{Composite Metric for Cyber Risk Prioritization:}
    Developing a vulnerability prioritization metric combining CVSS, CWSS, EPSS, and LEV.

    \item \textbf{Predictive Analytics for Russet Potato Selection:}
    Applying machine learning and feature selection to phenotypic data for potato breeding decisions.

    \item \textbf{Machine Learning for USEPA SPECIATE Database:}
    Applying clustering methods to identify similarities among chemical source profiles.

    \item \textbf{Miscellaneous Collaborative Research:}
    Collaborating on interdisciplinary research in graph learning, cybersecurity, environmental science, and applied AI.
\end{itemize}

\textbf{More Projects:} A complete list of my other projects is available \href{https://sites.google.com/site/siddhartha047/}{here}.

\end{rSection}

\begin{rSection}{Teaching and Mentoring} 

\textbf{Pacific Northwest National Laboratory (PNNL)}\hfill Aug 2025 -- May 2026\\
\textbf{Research Mentoring:} Mentored five interns over the past year on graph neural networks, circuit representation learning, and cybersecurity, resulting in three accepted publications.

\textbf{Purdue University, West Lafayette, IN, USA}\hfill Aug 2018 -- May 2019\\
Department of Computer Science\\
\textbf{Position:} Graduate Research Assistant.\\
\textbf{Courses:} Algorithms, Data Structures, etc.




\textbf{Bangladesh University of Engineering and Technology (BUET), Dhaka, Bangladesh}\hfill Aug 2014 -- Jul 2018\\
Department of Computer Science and Engineering\\
\textbf{Position:} Assistant Professor (2017-2018), Lecturer (2014-2017).\\
\textbf{Courses:} Machine Learning, Artificial Intelligence, Computer Graphics, Data Communication, and others (\href{https://sites.google.com/site/siddhartha047/courses}{List of Courses})

\end{rSection}
% \clearpage
\begin{rSection}{Skills}
\begin{tabularx}{\linewidth}{@{}>{\raggedright\arraybackslash}p{0.20\linewidth}X@{}}
\textbf{Research Areas:} & AI/ML, Graph neural networks, Graph Algorithms, Graph Sparsification, Natural Language Processing, Cybersecurity, Scientific Machine Learning.\\

\textbf{Languages:} & Python, C/C++, Java, MATLAB, Julia, C\#, SQL, JavaScript.\\

\textbf{Frameworks/Tools:} & PyTorch, PyTorch Geometric, DGL, TensorFlow, Keras, MPI, Git, Linux, MOEA Framework.\\

\textbf{LLM and AI Systems:} & Large language models, multimodal AI, agentic AI, retrieval-augmented generation (RAG), document intelligence, metadata extraction, georeferencing, interactive AI assistants.\\

\textbf{Systems and Applications:} & System design, hardware/software partitioning, scalable AI systems, large-scale document processing, vulnerability analysis, CVE--CWE--CAPEC mapping, cyber risk prioritization.\\

\textbf{Databases:} & Oracle, Microsoft SQL Server, MySQL, SQLite.\\
\end{tabularx}
\end{rSection}

\begin{rSection}{Academic Services} 
\textbf{Reviewer:} \href{https://icml.cc/Conferences/2025}{ICML'25}, \href{https://kdd2024.kdd.org/}{KDD'24}, \href{http://dsaa2022.dsaa.co/}{DSAA'23},
\href{https://www2021.thewebconf.org/}{The Web Conference (formerly WWW)}, 
\href{https://cis.ieee.org/publications/ieee-transactions-on-artificial-intelligence}{IEEE Transactions on Artificial Intelligence} and other venues.

\textbf{Organizing Member:} NSysS \href{https://cse.buet.ac.bd/nsyss/}{2015}, \href{http://cse.buet.ac.bd/nsyss2016/}{2016}, \href{https://cse.buet.ac.bd/nsyss2017/}{2017}.
\end{rSection}

\begin{rSection}{Awards and Achievements} 
\href{https://www.cs.purdue.edu/news/articles/2021/pothen_v2w-bert.html}{Best Application Paper Award, DSAA 2021}\\
\href{https://lh4.googleusercontent.com/RgzcFOh0gfXK4MTvMchVhGQ-X39kd13UFPGsZuHLkCcC9-CJNQ2dW_nfg0pIcuCSYbpbWlt94_z1dST95MYITwod_aW0tvTmpOLhRZqZs9eC73thuy5eA4YSDZLll9EK-w=w1280}{Best Paper Award, HST 2022}\\
Dean’s Honor List, BUET, 2011--2014\\
University Merit Scholarship, BUET
\end{rSection}


\begin{rSection}{References} 
\begin{table}[H]
\centering
\addtolength{\tabcolsep}{20pt}
\begin{tabular}{ll}
\begin{tabular}[c]{@{}l@{}}
\href{https://www.cs.purdue.edu/homes/apothen/}{\textbf{Alex Pothen}}\\
Professor,\\
Department of Computer Science,\\
Purdue University,\\
West Lafayette, IN, USA.\\
Email: \href{mailto:apothen@purdue.edu}{apothen@purdue.edu}
\end{tabular}
&
\begin{tabular}[c]{@{}l@{}}
\href{https://www.pnnl.gov/science/staff/staff_info.asp?staff_num=7452}{\textbf{Mahantesh Halappanavar}}\\
Scientist and Acting Group Leader,\\
Advanced Computing, Mathematics and Data,\\
Pacific Northwest National Laboratory,\\
Richland, WA, USA.\\
Email: \href{mailto:hala@pnnl.gov}{hala@pnnl.gov}
\end{tabular} \\\\

\begin{tabular}[c]{@{}l@{}}
\href{https://www.boisestate.edu/coen-cs/people/faculty/edoardo-serra/}{\textbf{Edoardo Serra}}\\
Associate Professor,\\
Department of Computer Science,\\
Boise State University,\\
Boise, ID, USA.\\
Email: \href{mailto:edoardoserra@boisestate.edu}{edoardoserra@boisestate.edu}
\end{tabular}
&
\begin{tabular}[c]{@{}l@{}}
\href{https://cse.buet.ac.bd/faculty/facdetail.php?id=mdmonirulislam}{\textbf{Md. Monirul Islam}}\\
Professor,\\
Department of Computer Science,\\
Bangladesh University of Engineering and Technology,\\
Dhaka, Bangladesh.\\
Email: \href{mailto:mmislam@cse.buet.ac.bd}{mmislam@cse.buet.ac.bd}
\end{tabular}
\end{tabular}
\end{table}
\end{rSection}

\end{document}